\documentclass[10pt, aspectratio=169]{beamer}
% Preámbulo modular para presentaciones Beamer (Modelación Estadística)
% Diseño y paleta institucional del Tecnológico de Monterrey

\documentclass[aspectratio=169,xcolor={dvipsnames,table}]{beamer}

\usepackage[utf8]{inputenc}
\usepackage[spanish,mexico]{babel}
\usepackage{amsmath,amssymb,amsfonts,mathtools}
\usepackage{graphicx}
\usepackage{listings}
\usepackage{booktabs}
\usepackage{tikz}
\usetikzlibrary{matrix,arrows,decorations.pathmorphing,shapes.geometric,positioning}

% Paleta Institucional Tecnológico de Monterrey
\definecolor{TecAzul}{HTML}{0039A6}       % Azul TEC: color primario institucional
\definecolor{TecAzulOscuro}{HTML}{1A2E51} % Azul marino oscuro
\definecolor{TecRojo}{HTML}{EC2661}       % Rojo de acento (paleta miTec)
\definecolor{TecGris}{HTML}{646464}       % Gris neutro para comentarios y subtítulos
\definecolor{TecFondoBloque}{HTML}{F4F6F9}% Fondo suave para bloques

% Configuración de Tema Beamer
\usetheme{Boadilla}
\usecolortheme[named=TecAzulOscuro]{structure}
\setbeamercolor{alerted text}{fg=TecRojo}
\setbeamercolor{palette primary}{bg=TecAzulOscuro,fg=white}
\setbeamercolor{palette secondary}{bg=TecAzul,fg=white}
\setbeamercolor{palette tertiary}{bg=TecAzulOscuro!80!black,fg=white}
\setbeamercolor{title}{fg=TecAzulOscuro,bg=TecFondoBloque}
\setbeamercolor{frametitle}{fg=TecAzulOscuro,bg=TecFondoBloque}

% Bloques personalizados elegantes
\setbeamercolor{block title}{bg=TecAzulOscuro,fg=white}
\setbeamercolor{block body}{bg=TecFondoBloque,fg=black}
\setbeamercolor{block title alerted}{bg=TecRojo,fg=white}
\setbeamercolor{block body alerted}{bg=TecRojo!8,fg=black}
\setbeamercolor{block title example}{bg=TecAzul,fg=white}
\setbeamercolor{block body example}{bg=TecAzul!8,fg=black}

% Redondear bordes y añadir sombra sutil
\setbeamertemplate{blocks}[rounded][shadow=true]
\setbeamertemplate{navigation symbols}{}

% Pie de página limpio con número de diapositiva
\setbeamertemplate{footline}{
  \leavevmode%
  \hbox{%
  \begin{beamercolorbox}[wd=.7\paperwidth,ht=2.25ex,dp=1ex,left]{author in head/foot}%
    \hspace*{2ex}\insertshortauthor~~(\insertshortinstitute) \hfill \insertshorttitle
  \end{beamercolorbox}%
  \begin{beamercolorbox}[wd=.3\paperwidth,ht=2.25ex,dp=1ex,right]{date in head/foot}%
    \insertshortdate{}\hspace*{2em}
    \insertframenumber{} / \inserttotalframenumber\hspace*{2ex}
  \end{beamercolorbox}}%
  \vskip0pt%
}

% Configuración de Listings de Python para visualización pedagógica de código
\lstset{
    language=Python,
    basicstyle=\ttfamily\footnotesize,
    keywordstyle=\color{TecAzulOscuro}\bfseries,
    stringstyle=\color{TecRojo},
    commentstyle=\color{TecGris}\itshape,
    showstringspaces=false,
    breaklines=true,
    frame=lines,
    rulecolor=\color{TecAzulOscuro!30},
    backgroundcolor=\color{TecFondoBloque},
    aboveskip=0.5em,
    belowskip=0.5em,
    literate={á}{{\'a}}1 {é}{{\'e}}1 {í}{{\'i}}1 {ó}{{\'o}}1 {ú}{{\'u}}1
             {Á}{{\'A}}1 {É}{{\'E}}1 {Í}{{\'I}}1 {Ó}{{\'O}}1 {Ú}{{\'U}}1
             {ñ}{{\~n}}1 {Ñ}{{\~N}}1 {¿}{{?`}}1 {¡}{{!`}}1
}

% Metadatos institucionales por defecto
\title[Modelación Estadística]{Modelación Estadística para Ciencia de Datos e Ingeniería}
\author[J. Castillo Colmenares]{Juliho David Castillo Colmenares}
\institute[Tec de Monterrey]{Tecnológico de Monterrey \\ Escuela de Ingeniería y Ciencias}
\date{\today}

% Comandos y macros estadísticas compartidas para las presentaciones Beamer (Metropolis)

% Conjuntos numéricos básicos
\newcommand{\R}{\mathbb{R}}
\newcommand{\N}{\mathbb{N}}
\newcommand{\Q}{\mathbb{Q}}
\newcommand{\Z}{\mathbb{Z}}
\newcommand{\set}[1]{\left\{ #1 \right\}}
\newcommand{\abs}[1]{\left|#1\right|}
\newcommand{\norm}[1]{\left\| #1 \right\|}

% Comandos estadísticos y probabilísticos del curso
\newcommand{\Var}{\operatorname{Var}}
\newcommand{\cov}{\operatorname{Cov}}
\newcommand{\corr}{\rho}
\newcommand{\comb}[2]{\begin{pmatrix} #1 \\ #2 \end{pmatrix}}
\newcommand{\s}{\sigma}
\newcommand{\particion}{\mathfrak{P}}
\newcommand{\card}[1]{\#\left( #1 \right)}
\newcommand{\seq}[3]{\set{#1}_{#2}^{#3}}
\newcommand{\E}{\mathbb{E}}
\newcommand{\Prob}{\mathbb{P}}

% Entornos pedagógicos y cajas en español académico natural
\newenvironment{discusion}{%
  \begin{alertblock}{Pregunta de Discusión}%
}{%
  \end{alertblock}%
}

\newenvironment{reflexion}{%
  \begin{alertblock}{Reflexión para el Aula}%
}{%
  \end{alertblock}%
}

\newenvironment{paradoja}{%
  \begin{alertblock}{Paradoja de Exploración}%
}{%
  \end{alertblock}%
}

\newenvironment{motivacion}{%
  \begin{exampleblock}{Motivación: Ciencia de Datos en la Práctica}%
}{%
  \end{exampleblock}%
}

\newenvironment{retoclase}{%
  \begin{block}{Reto Rápido en Equipo}%
}{%
  \end{block}%
}


\title[Estadísticos y Varianza Insesgada]{Sección 05.01: Muestreo Aleatorio, Media y Varianza Muestral Insesgada}
\subtitle{Estadísticos, Corrección de Bessel y Corrección por Población Finita}
\author[J. Castillo Colmenares]{Juliho Castillo Colmenares}
\institute{Tecnológico de Monterrey}
\date{\vspace{-1.2cm}}

\begin{document}

% Slide 1: Portada
\begin{frame}[plain]
  \vspace{-0.8cm}
  \titlepage
\end{frame}

% Slide 2: Hoja de Ruta
\begin{frame}{Hoja de Ruta --- Capítulo 05: Distribuciones de Muestreo}
  \footnotesize
  ¿Dónde nos encontramos en el desarrollo modular del capítulo? \pause
  \begin{itemize}\itemsep=0.08em
    \item \textbf{05.01} Muestreo Aleatorio Simple, Media y Varianza Muestral Insesgada \emph{(Hoy)} \pause
    \item \textbf{05.02} Teorema del Límite Central Asintótico \pause
    \item \textbf{05.03} Distribución Chi-Cuadrada y Varianza Muestral \pause
    \item \textbf{05.04} Distribución $t$ de Student y Muestras Pequeñas \pause
    \item \textbf{05.05} Distribución $F$ de Fisher-Snedecor
  \end{itemize}
  \vspace{-0.05cm}
  \begin{block}{Objetivo de la Sesión}
    Formalizar el concepto de \emph{estadístico} como variable aleatoria, dominar la distribución muestral de la media ($E(\bar{X})=\mu$, $\Var(\bar{X})=\sigma^2/n$), y demostrar por qué la varianza muestral $S^2$ debe dividirse entre $n-1$ (corrección de Bessel) para ser un estimador insesgado de $\sigma^2$.
  \end{block}
\end{frame}

% Slide 3: Motivación
\begin{frame}{De las Distribuciones a los Estadísticos}
  \footnotesize
  \begin{motivacion}
    Hasta ahora hemos estudiado variables aleatorias individuales. En inferencia estadística, trabajamos con \emph{muestras}: colecciones de observaciones $X_1, \ldots, X_n$. Cualquier función de la muestra --- como la media o la varianza muestral --- es en sí misma una variable aleatoria, con su propia distribución de probabilidad.
  \end{motivacion}

  \vspace{0.15cm}
  \pause
  \begin{itemize}\itemsep=0.1cm
    \item \textbf{Estadístico:} función de la muestra que no depende de parámetros desconocidos ($\bar{X}$, $S^2$, $\hat{p}$). \pause
    \item \textbf{Distribución muestral:} la distribución de probabilidad de un estadístico. \pause
    \item \textbf{Pregunta central:} ¿cómo se comporta $\bar{X}$ y $S^2$ al repetir el muestreo?
  \end{itemize}
\end{frame}

% Slide 4: Estadístico y Distribución Muestral de la Media
\begin{frame}{Estadísticos y la Distribución Muestral de la Media}
  \footnotesize
  Sea $X_1, \ldots, X_n$ una muestra aleatoria simple (iid) de una población con media $\mu$ y varianza $\sigma^2$. \pause

  \vspace{0.1cm}
  \begin{block}{Media y Varianza de $\bar{X}$}
    $$E(\bar{X}) = \mu, \qquad \Var(\bar{X}) = \dfrac{\sigma^2}{n}, \qquad \text{DE}(\bar{X}) = \dfrac{\sigma}{\sqrt{n}}$$
    $\bar{X}$ es un estimador \emph{insesgado} de $\mu$, y su dispersión disminuye con $n$: al cuadruplicar $n$, el error estándar se reduce a la mitad.
  \end{block}

  \vspace{0.1cm}
  \pause
  \begin{alertblock}{Consistencia}
    \begin{itemize}\itemsep=0.05cm
      \item Como $\Var(\bar{X}) = \sigma^2/n \to 0$ cuando $n\to\infty$, y $E(\bar{X})=\mu$ para todo $n$, $\bar{X}$ converge en probabilidad a $\mu$ (desigualdad de Chebyshev).
      \item Decimos que $\bar{X}$ es un estimador \emph{consistente} de $\mu$.
    \end{itemize}
  \end{alertblock}
\end{frame}

% Slide 5: Varianza Muestral Insesgada
\begin{frame}{Varianza Muestral Insesgada: Corrección de Bessel}
  \footnotesize
  Definimos la varianza muestral con divisor $n-1$: \pause

  \vspace{0.1cm}
  \begin{block}{Definición y Teorema de Insesgadez}
    $$S^2 = \dfrac{1}{n-1}\sum_{i=1}^n (X_i-\bar{X})^2, \qquad E(S^2) = \sigma^2$$
    El divisor $n-1$ no es arbitrario: es exactamente el que hace a $S^2$ insesgado.
  \end{block}

  \vspace{0.1cm}
  \pause
  \begin{alertblock}{¿Por qué no dividir entre $n$?}
    \begin{itemize}\itemsep=0.05cm
      \item El estimador ``natural'' $\hat{\sigma}^2 = \frac{1}{n}\sum(X_i-\bar{X})^2$ tiene $E(\hat{\sigma}^2) = \frac{n-1}{n}\sigma^2 < \sigma^2$: subestima sistemáticamente.
      \item $\bar{X}$ está, por construcción, más cerca de cada $X_i$ que la media poblacional desconocida $\mu$; dividir entre $n-1$ (``perder un grado de libertad'') corrige ese sesgo.
    \end{itemize}
  \end{alertblock}
\end{frame}

% Slide 6: Demostración de la Insesgadez
\begin{frame}{Demostración: $E(S^2) = \sigma^2$}
  \footnotesize
  Escribimos $(X_i-\bar{X}) = (X_i-\mu) - (\bar{X}-\mu)$ y expandimos: \pause

  \vspace{0.1cm}
  \begin{block}{Desarrollo}
    $$\sum_{i=1}^n (X_i-\bar{X})^2 = \sum_{i=1}^n (X_i-\mu)^2 - n(\bar{X}-\mu)^2$$
    usando que $\sum(X_i-\mu) = n(\bar{X}-\mu)$ para cancelar el término cruzado.
  \end{block}

  \vspace{0.1cm}
  \pause
  \begin{alertblock}{Tomando Esperanza}
    Con $E[(X_i-\mu)^2]=\sigma^2$ y $E[(\bar{X}-\mu)^2]=\Var(\bar{X})=\sigma^2/n$:
    $$E\left[\sum(X_i-\bar{X})^2\right] = n\sigma^2 - n\cdot\frac{\sigma^2}{n} = (n-1)\sigma^2 \implies E(S^2)=\sigma^2.$$
  \end{alertblock}
\end{frame}

% Slide 7: Corrección por Población Finita
\begin{frame}{Muestreo sin Reemplazo: Corrección por Población Finita (FPC)}
  \footnotesize
  Al muestrear \emph{sin reemplazo} de una población finita de tamaño $N$, la fórmula de $\Var(\bar{X})$ se ajusta: \pause

  \vspace{0.1cm}
  \begin{block}{Factor de Corrección por Población Finita}
    $$\Var(\bar{X}) = \dfrac{\sigma^2}{n}\cdot\underbrace{\dfrac{N-n}{N-1}}_{\text{FPC}}$$
    El factor FPC es siempre $< 1$ cuando $n>1$, reduciendo la varianza respecto al caso de población infinita.
  \end{block}

  \vspace{0.1cm}
  \pause
  \begin{alertblock}{Interpretación}
    \begin{itemize}\itemsep=0.05cm
      \item Cuando $N\to\infty$ (o la fracción muestreada $n/N \to 0$), el factor FPC $\to 1$ y ambas fórmulas coinciden.
      \item Relevante en auditorías, encuestas electorales y control de calidad sobre lotes finitos.
    \end{itemize}
  \end{alertblock}
\end{frame}

% Slide 8: Aplicaciones
\begin{frame}{Aplicaciones en Inferencia y Ciencia de Datos}
  \footnotesize
  Los estadísticos $\bar{X}$ y $S^2$ son la base de toda la inferencia estadística: \pause

  \vspace{0.15cm}
  \begin{itemize}\itemsep=0.12cm
    \item \textbf{Intervalos de Confianza:} $S^2$ insesgado permite construir errores estándar confiables para $\mu$. \pause
    \item \textbf{Pruebas de Hipótesis:} el estadístico $t = (\bar{X}-\mu_0)/(S/\sqrt{n})$ (Capítulo 05.04) depende directamente de $S^2$. \pause
    \item \textbf{Auditoría y Control de Calidad:} FPC ajusta la incertidumbre al muestrear lotes finitos sin reemplazo. \pause
    \item \textbf{Machine Learning:} varianza insesgada en validación cruzada y estimación de error de generalización.
  \end{itemize}
\end{frame}

% Slide 9: Lab Python (Bloque 1)
\begin{frame}[fragile]{Laboratorio en Python: Varianza Muestral Insesgada}
  \scriptsize
  Verificación Monte Carlo de $E(S^2)=\sigma^2$ vs. el estimador sesgado, y la fórmula abreviada (\texttt{05.01\_sample\_statistics.py}):
  {\lstset{basicstyle=\fontsize{4pt}{4.8pt}\selectfont\ttfamily,aboveskip=0.1em,belowskip=0.1em}
  \lstinputlisting[firstline=17, lastline=43]{../../code/05_distribuciones_muestreo/05.01_sample_statistics.py}}
\end{frame}

% Slide 10: Lab Python (Bloque 2)
\begin{frame}[fragile]{Laboratorio en Python: Distribución Muestral de la Media}
  \scriptsize
  Verificación de $E(\bar{X})=\mu$ y $\Var(\bar{X})=\sigma^2/n$ para distintos tamaños de muestra:
  {\lstset{basicstyle=\fontsize{4pt}{4.8pt}\selectfont\ttfamily,aboveskip=0.1em,belowskip=0.1em}
  \lstinputlisting[firstline=46, lastline=68]{../../code/05_distribuciones_muestreo/05.01_sample_statistics.py}}
\end{frame}

% Slide 11: Lab Python (Bloque 3)
\begin{frame}[fragile]{Laboratorio en Python: Corrección por Población Finita}
  \scriptsize
  Comparación de la varianza ingenua vs. la FPC, con verificación empírica por muestreo sin reemplazo:
  {\lstset{basicstyle=\fontsize{4pt}{4.8pt}\selectfont\ttfamily,aboveskip=0.1em,belowskip=0.1em}
  \lstinputlisting[firstline=71, lastline=102]{../../code/05_distribuciones_muestreo/05.01_sample_statistics.py}}
\end{frame}

% Slide 12: Salida Terminal
\begin{frame}[fragile]{Laboratorio en Python: Salida en Terminal}
  \scriptsize
  Salida estándar generada al ejecutar el script del laboratorio en Python:
  \vspace{0.02cm}
  \begin{block}{Salida Estándar en Consola}
    \fontsize{4pt}{5pt}\selectfont
    \begin{verbatim}
=== Block 1: Unbiased Sample Variance (Bessel's Correction) ===
Population: N(50.0, 100.0), n=5, trials=200000
  E(S^2) [unbiased]: 100.1676 (theoretical: 100.0000)
  E(sigma_hat^2) [biased]: 80.1341 (theoretical: 80.0000)
  Worked example {8,12,10,14,6}: mean=10.0000, S^2=10.0000

=== Block 2: Sampling Distribution of the Mean ===
  n=25:  Var(Xbar)=63.98 (theo 64.00) | n=100: 16.04 (theo 16.00) | n=400: 4.00 (theo 4.00)

=== Block 3: Finite Population Correction (FPC) ===
N=200, n=20: naive Var=34.83, FPC factor=0.9045, FPC Var=31.50
Empirical Var(Xbar) [without replacement]: 31.61
    \end{verbatim}
  \end{block}
\end{frame}

% Slide 17: Cierre
\begin{frame}{Cierre de Sección: Estadísticos y Varianza Insesgada}
  \begin{reflexion}
    Todo estadístico calculado de una muestra --- la media, la varianza, la proporción --- es en sí mismo una variable aleatoria con su propia distribución muestral. Formalizar la insesgadez de $\bar{X}$ y $S^2$ es el primer paso para construir intervalos de confianza y pruebas de hipótesis rigurosas en las secciones siguientes.
  \end{reflexion}

  \vspace{0.2cm}
  \pause
  \begin{block}{Lecciones Clave}
    \begin{itemize}\itemsep=0.08cm
      \item \textbf{Distribución muestral de $\bar{X}$:} $E(\bar{X})=\mu$, $\Var(\bar{X})=\sigma^2/n$; estimador insesgado y consistente.
      \item \textbf{Corrección de Bessel:} $S^2$ divide entre $n-1$, no $n$, para ser insesgado: $E(S^2)=\sigma^2$.
      \item \textbf{FPC:} al muestrear sin reemplazo de una población finita, $\Var(\bar{X})$ se multiplica por $(N-n)/(N-1) < 1$.
    \end{itemize}
  \end{block}
\end{frame}

% Slide 18: Perspectiva Modular
\begin{frame}{Perspectiva Modular --- El Próximo Reto}
  \scriptsize
  \begin{retoclase}
    Con la media y varianza muestral formalizadas, la Sección 05.02 profundizará en el Teorema del Límite Central de forma asintótica y rigurosa, y las Secciones 05.03--05.05 introducirán las distribuciones Chi-cuadrada, $t$ de Student y $F$ de Fisher-Snedecor: las herramientas exactas que sustituyen a la aproximación Normal cuando la muestra es pequeña o se desconoce $\sigma$.
  \end{retoclase}

  \vspace{0.08cm}
  \pause
  \begin{alertblock}{Preparación Recomendada}
    \begin{itemize}\itemsep=0.06cm
      \item Repasa la conexión Gamma--Chi-cuadrada (Sección 04.07): $\chi^2_\nu = \text{Gamma}(\nu/2, 2)$.
      \item Reflexiona sobre por qué $(n-1)S^2/\sigma^2 \sim \chi^2_{n-1}$: la varianza muestral insesgada de hoy es la pieza central de esa distribución.
    \end{itemize}
  \end{alertblock}
\end{frame}

\end{document}
